\documentclass{article}

\begin{document}
\subsection*{Problem 4}

Each of four home mortgages is classified as fixed rate ($F$) or variable rate ($V$). An outcome is an ordered
4-tuple of $F$'s and $V$'s.

\begin{enumerate}
\item[(a)] The sample space has $2^4=16$ outcomes:
\[
\mathcal{S}=\left\{\begin{array}{l}
FFFF,FFFV,FFVF,FFVV,FVFF,FVFV,FVVF,FVVV,\\
VFFF,VFFV,VFVF,VFVV,VVFF,VVFV,VVVF,VVVV
\end{array}\right\}.
\]

\item[(b)] Event $B=$ ``exactly three are fixed'' $\iff$ exactly one $V$:
\[
B=\{VFFF,FVFF,FFVF,FFFV\}.
\]

\item[(c)] Event $C=$ ``all four are the same type'':
\[
C=\{FFFF,\,VVVV\}.
\]

\item[(d)] Event $D=$ ``at most one is variable'' $\iff$ 0 or 1 $V$:
\[
D=\{FFFF\}\cup\{VFFF,FVFF,FFVF,FFFV\}=\{FFFF,VFFF,FVFF,FFVF,FFFV\}.
\]

\item[(e)] Union and intersection of events in (c) and (d):
\[
C\cup D=\{FFFF,VVVV\}\cup\{FFFF,VFFF,FVFF,FFVF,FFFV\}
\]
\[
      =\{FFFF,VFFF,FVFF,FFVF,FFFV,VVVV\}
\]
\[
C\cap D=\{FFFF,VVVV\}\cap\{FFFF,VFFF,FVFF,FFVF,FFFV\}=\{FFFF\}.
\]

\item[(f)] Union and intersection of events in (b) and (c):
\[
B\cup C=\{VFFF,FVFF,FFVF,FFFV\}\cup\{FFFF,VVVV\}
\]
\[
      =\{VFFF,FVFF,FFVF,FFFV,FFFF,VVVV\},
\]
\[
B\cap C=\{VFFF,FVFF,FFVF,FFFV\}\cap\{FFFF,VVVV\}=nothing.
\]
\end{enumerate}

\subsection*{Problem 6}

Books $1,2$ are first printings; books $3,4,5$ are second printings. The student examines books in random
order and stops when the first second printing is selected. Thus, each outcome is a sequence that ends with
the first occurrence of an element of $\{3,4,5\}$, and any preceding terms (if any) must be from $\{1,2\}$
with no repetition.

\begin{enumerate}
\item[(a)] The possible outcomes are:
\[
\mathcal{S}=
\{3,4,5,\;
13,14,15,\;
23,24,25,\;
213,214,215,\;
123,124,125\}.
\]

\item[(b)] Let $A=$ ``exactly one book is examined.'' This happens iff the first book examined is a second
    printing:
\[
A=\{3,4,5\}.
\]

\item[(c)] Let $B=$ ``book 5 is the one selected'' (i.e., the stopping book is 5). Then the outcome must end
    in 5, with any preceding terms drawn from $\{1,2\}$:
\[
B=\{5,\,15,\,25,\,125,\,215\}.
\]

\item[(d)] Let $C=$ ``book 1 is not examined.'' So the outcome cannot contain 1. The possible outcomes are
    those consisting of either a single second-printing book, or starting with 2 (possibly) and then a
    second-printing book:
\[
C=\{3,4,5,\,23,24,25\}.
\]
\end{enumerate}

\subsection*{Problem 12}

Let $A$ be the event the selected student has a Visa card and $B$ the event the student has a MasterCard.
Given
\[
P(A)=0.6,\qquad P(B)=0.4.
\]

\begin{enumerate}
\item[(a)] Could $P(A\cap B)=0.5$?
\[
P(A\cap B)\le \min\{P(A),P(B)\}=\min\{0.6,0.4\}=0.4.
\]
Thus $P(A\cap B)$ cannot be $0.5$.

\item[(b)] Now suppose $P(A\cap B)=0.3$. Probability of at least one card is
\[
P(A\cup B)=P(A)+P(B)-P(A\cap B)=0.6+0.4-0.3=0.7.
\]

\item[(c)] Probability of neither type is
\[
P\big((A\cup B)^c\big)=1-P(A\cup B)=1-0.7=0.3.
\]

\item[(d)] ``Visa but not MasterCard'' is the event $A\cap B^c$ (equivalently $A\setminus B$). Its
    probability:
\[
P(A\cap B^c)=P(A)-P(A\cap B)=0.6-0.3=0.3.
\]

\item[(e)] ``Exactly one'' of the two types is the symmetric difference
\[
(A\cap B^c)\ \cup\ (A^c\cap B).
\]
Compute
\[
P(A^c\cap B)=P(B)-P(A\cap B)=0.4-0.3=0.1,
\]
and these two events are disjoint, so
\[
P(\text{exactly one})=P(A\cap B^c)+P(A^c\cap B)=0.3+0.1=0.4.
\]
\end{enumerate}

\subsection*{Problem 18}

A wallet has $5$ ten-dollar bills, $4$ five-dollar bills, and $6$ one-dollar bills, for a total of
\[
5+4+6=15 \text{ bills.}
\]
Bills are drawn without replacement in random order. Let $A$ be the event that \emph{at least two bills must
be selected to obtain the first \$10 bill}. This means the first bill drawn is \emph{not} a \$10 bill.

There are $4+6=10$ non-\$10 bills, so
\[
P(A)=P(\text{first bill is not \$10})=\frac{10}{15}=\frac{2}{3}\approx 0.6667.
\]


\subsection*{Problem 24}

Assume $A\subseteq B$. Then $A$ and $B\cap A^{c}$ are disjoint and
\[
B = A \,\cup\, (B\cap A^{c}).
\]
By additivity on disjoint events,
\[
P(B)=P(A)+P(B\cap A^{c}) \ge P(A)
\]
since $P(B\cap A^{c})\ge 0$. Therefore,
\[
\boxed{\,P(A)\le P(B)\,}.
\]

For general events $A$ and $B$, we have the containments
\[
A\cap B \subseteq A \subseteq A\cup B.
\]
Applying the result just proved to each inclusion gives
\[
\boxed{\,P(A\cap B)\le P(A)\le P(A\cup B)\,}.
\]

\subsection*{Problem 14}

Let $C=$ event an adult regularly consumes coffee, and $S=$ event an adult regularly consumes carbonated soda.
Given
\[
P(C)=0.55,\quad P(S)=0.45,\quad P(C\cup S)=0.70.
\]

\begin{enumerate}
\item[(a)] Use the addition rule:
\begin{align*}
P(C\cap S)
&= P(C)+P(S)-P(C\cup S) \\
&= 0.55+0.45-0.70 \\
&= 0.30.
\end{align*}

\item[(b)] Interpret ``doesn't regularly consume at least one of these two products'' as ``consumes neither
    coffee nor soda,'' i.e.\ $(C\cup S)^c$.
\[
P\big((C\cup S)^c\big)=1-P(C\cup S)=1-0.70=0.30.
\]
\end{enumerate}

\subsection*{Problem 26}

Let $A_i$ denote the event the system has a defect of type $i$ ($i=1,2,3$). Given:
\[
P(A_1)=0.12,\quad P(A_2)=0.07,\quad P(A_3)=0.05,
\]
\[
P(A_1\cup A_2)=0.13,\quad P(A_1\cup A_3)=0.14,\quad P(A_2\cup A_3)=0.10,
\]
\[
P(A_1\cap A_2\cap A_3)=0.01.
\]

First find pairwise intersections using $P(A\cup B)=P(A)+P(B)-P(A\cap B)$:
\begin{align*}
P(A_1\cap A_2) &= P(A_1)+P(A_2)-P(A_1\cup A_2)=0.12+0.07-0.13=0.06,\\
P(A_1\cap A_3) &= 0.12+0.05-0.14=0.03,\\
P(A_2\cap A_3) &= 0.07+0.05-0.10=0.02.
\end{align*}

\begin{enumerate}
\item[(a)] Probability of no type 1 defect:
\[
P(A_1^c)=1-P(A_1)=1-0.12=0.88.
\]

\item[(b)] Probability of both type 1 and type 2 defects:
\[
P(A_1\cap A_2)=0.06.
\]

\item[(c)] Probability of type 1 and type 2 but not type 3:
\[
P(A_1\cap A_2\cap A_3^c)=P(A_1\cap A_2)-P(A_1\cap A_2\cap A_3)=0.06-0.01=0.05.
\]

\item[(d)] ``At most two'' defects means \emph{not all three} defects:
\[
P(\text{at most two}) = 1 - P(A_1\cap A_2\cap A_3)=1-0.01=0.99.
\]
\end{enumerate}

\subsection*{Problem 30}

There are $8$ bottles of zinfandel (Z), $10$ of merlot (M), and $12$ of cabernet (C), for a total of $30$
distinct bottles.

\begin{enumerate}
\item[(a)] Serve $3$ bottles of zinfandel with serving order important:
\[
{}_{8}P_{3}=\frac{8!}{(8-3)!}=8\cdot 7\cdot 6=336.
\]

\item[(b)] Randomly select $6$ bottles from $30$ (order not important):
\[
\binom{30}{6}=\frac{30!}{6!\,24!}=593{,}775.
\]

\item[(c)] Select exactly $2$ of each variety (order not important):
\[
\binom{8}{2}\binom{10}{2}\binom{12}{2}
=28\cdot 45\cdot 66
=83{,}160.
\]

\item[(d)] Probability of getting two of each variety when $6$ are selected:
\[
P(\text{2 of each})=\frac{\binom{8}{2}\binom{10}{2}\binom{12}{2}}{\binom{30}{6}}
=\frac{83{,}160}{593{,}775}
=\frac{1848}{13195}
\approx 0.1400.
\]

\item[(e)] Probability all $6$ selected are the same variety:
\[
\#(\text{all same})=\binom{8}{6}+\binom{10}{6}+\binom{12}{6}
=28+210+924=1162,
\]
\[
P(\text{all same})=\frac{1162}{\binom{30}{6}}
=\frac{1162}{593{,}775}
=\frac{166}{84{,}825}
\approx 0.001956.
\]
\end{enumerate}

\subsection*{Problem 34}

There are $25$ failed keyboards: $6$ with electrical defects (E) and $19$ with mechanical defects (M).

\begin{enumerate}
\item[(a)] Number of ways to select $5$ keyboards (order irrelevant):
\[
\binom{25}{5}=\frac{25!}{5!\,20!}=53{,}130.
\]

\item[(b)] Exactly two with an electrical defect $\Rightarrow$ choose $2$ from the $6$ electrical and $3$ from
    the $19$ mechanical:
\[
\binom{6}{2}\binom{19}{3}=15\cdot 969=14{,}535.
\]

\item[(c)] For a random sample of $5$, probability at least $4$ have a mechanical defect:
\[
P(\text{at least }4\text{ M})=\frac{\binom{19}{4}\binom{6}{1}+\binom{19}{5}\binom{6}{0}}{\binom{25}{5}}.
\]
Compute:
\[
\binom{19}{4}\binom{6}{1}+\binom{19}{5}=6\cdot 3876+11628=34884,
\qquad \binom{25}{5}=53130,
\]
so
\[
P=\frac{34884}{53130}=\frac{5814}{8855}\approx 0.6566.
\]
\end{enumerate}

\subsection*{Problem 36}

There are $3$ votes for $A$ and $2$ votes for $B$. Each ordering of the $5$ slips is equally likely, so the
sample space size is
\[
|\mathcal S|=\binom{5}{2}=10
\]
(choose the $2$ positions occupied by $B$).

We want the probability that $A$ is \emph{strictly ahead} of $B$ after each slip is tallied.
By the (Bertrand) Ballot Theorem, for $p>q$ the number of sequences with $p$ $A$'s and $q$ $B$'s such that $A$
is always ahead is
\[
\frac{p-q}{p+q}\binom{p+q}{q}.
\]
Here $p=3$, $q=2$, so the number of favorable sequences is
\[
\frac{3-2}{3+2}\binom{5}{2}=\frac{1}{5}\cdot 10 = 2.
\]
(Indeed, the favorable orderings are $AAABB$ and $AABAB$.)

Therefore,
\[
P(\text{$A$ remains ahead throughout})=\frac{2}{10}=\frac{1}{5}=0.2.
\]


\subsection*{Problem 50}

The tables give joint probabilities for \emph{(sleeve length, size, pattern)}. Let Sh = short-sleeved, Lo =
long-sleeved; sizes $S,M,L$; patterns Pl (plaid), Pr (print), St (stripe).

Short-sleeved probabilities:
\[
\begin{array}{c|ccc}
 & \text{Pl} & \text{Pr} & \text{St}\\ \hline
S & 0.04 & 0.02 & 0.05\\
M & 0.08 & 0.07 & 0.12\\
L & 0.03 & 0.07 & 0.08
\end{array}
\]
Long-sleeved probabilities:
\[
\begin{array}{c|ccc}
 & \text{Pl} & \text{Pr} & \text{St}\\ \hline
S & 0.03 & 0.02 & 0.03\\
M & 0.10 & 0.05 & 0.07\\
L & 0.04 & 0.02 & 0.08
\end{array}
\]

\begin{enumerate}
\item[(a)]
\[
P(M\cap \text{Lo}\cap \text{Pr}) = 0.05.
\]

\item[(b)] Medium print (either sleeve length):
\[
P(M\cap \text{Pr}) = P(M\cap \text{Sh}\cap \text{Pr}) + P(M\cap \text{Lo}\cap \text{Pr})
=0.07+0.05=0.12.
\]

\item[(c)]
\[
P(\text{Sh}) = (0.04+0.02+0.05)+(0.08+0.07+0.12)+(0.03+0.07+0.08)=0.56,
\]
\[
P(\text{Lo}) = 1-P(\text{Sh}) = 1-0.56=0.44.
\]

\item[(d)] Size medium:
\[
P(M)=(0.08+0.07+0.12)+(0.10+0.05+0.07)=0.27+0.22=0.49.
\]
Pattern print:
\[
P(\text{Pr})=(0.02+0.07+0.07)+(0.02+0.05+0.02)=0.16+0.09=0.25.
\]

\item[(e)] Given short-sleeved plaid:
\[
P(M\mid \text{Sh}\cap \text{Pl})=\frac{P(M\cap \text{Sh}\cap \text{Pl})}{P(\text{Sh}\cap \text{Pl})}
=\frac{0.08}{0.04+0.08+0.03}
=\frac{0.08}{0.15}
=\frac{8}{15}\approx 0.5333.
\]

\item[(f)] Given medium plaid:
\[
P(\text{Sh}\mid M\cap \text{Pl})=\frac{P(\text{Sh}\cap M\cap \text{Pl})}{P(M\cap \text{Pl})}
=\frac{0.08}{0.08+0.10}
=\frac{0.08}{0.18}
=\frac{4}{9}\approx 0.4444,
\]
\[
P(\text{Lo}\mid M\cap \text{Pl})=\frac{0.10}{0.08+0.10}
=\frac{0.10}{0.18}
=\frac{5}{9}\approx 0.5556.
\]
\end{enumerate}

\subsection*{Problem 56}

Let $A$ and $B$ be events with $P(B)>0$. By definition of conditional probability,
\[
P(A\mid B)=\frac{P(A\cap B)}{P(B)},\qquad
P(A^c\mid B)=\frac{P(A^c\cap B)}{P(B)}.
\]
Note that $A\cap B$ and $A^c\cap B$ are disjoint and their union is $B$:
\[
(A\cap B)\cup(A^c\cap B)=B,\qquad (A\cap B)\cap(A^c\cap B)=\varnothing.
\]
Therefore,
\[
P(A\cap B)+P(A^c\cap B)=P\big((A\cap B)\cup(A^c\cap B)\big)=P(B).
\]
Divide both sides by $P(B)>0$:
\[
P(A\mid B)+P(A^c\mid B)=\frac{P(A\cap B)}{P(B)}+\frac{P(A^c\cap B)}{P(B)}
=\frac{P(B)}{P(B)}=1.
\]


\subsection*{Problem 58}

Let $A,B,C$ be events with $P(C)>0$. By definition of conditional probability,
\[
P(A\cup B\mid C)=\frac{P\big((A\cup B)\cap C\big)}{P(C)}.
\]
Use distributivity and the addition rule:
\[
(A\cup B)\cap C=(A\cap C)\cup(B\cap C),
\]
and
\[
P\big((A\cap C)\cup(B\cap C)\big)=P(A\cap C)+P(B\cap C)-P(A\cap B\cap C).
\]
Therefore,
\begin{align*}
P(A\cup B\mid C)
&=\frac{P(A\cap C)+P(B\cap C)-P(A\cap B\cap C)}{P(C)}\\
&=\frac{P(A\cap C)}{P(C)}+\frac{P(B\cap C)}{P(C)}-\frac{P(A\cap B\cap C)}{P(C)}\\
&=P(A\mid C)+P(B\mid C)-P(A\cap B\mid C).
\end{align*}


\subsection*{Problem 60}

Let $D$ = event the aircraft is discovered, and $L$ = event it has an emergency locator.
Given:
\[
P(D)=0.70 \;\Rightarrow\; P(D^c)=0.30,\qquad P(L\mid D)=0.60,
\]
\[
P(L^c\mid D^c)=0.90 \;\Rightarrow\; P(L\mid D^c)=0.10.
\]
First compute
\[
P(L)=P(L\mid D)P(D)+P(L\mid D^c)P(D^c)
=0.60(0.70)+0.10(0.30)=0.42+0.03=0.45.
\]

\begin{enumerate}
\item[(a)] If it has an emergency locator, the probability it will \emph{not} be discovered:
\[
P(D^c\mid L)=\frac{P(L\mid D^c)P(D^c)}{P(L)}
=\frac{0.10(0.30)}{0.45}=\frac{0.03}{0.45}=\frac{1}{15}\approx 0.0667.
\]

\item[(b)] If it does \emph{not} have an emergency locator, the probability it \emph{will} be discovered:

First,
\[
P(L^c)=1-P(L)=1-0.45=0.55,\qquad P(L^c\mid D)=1-0.60=0.40.
\]
Then
\[
P(D\mid L^c)=\frac{P(L^c\mid D)P(D)}{P(L^c)}
=\frac{0.40(0.70)}{0.55}=\frac{0.28}{0.55}=\frac{28}{55}\approx 0.5091.
\]
\end{enumerate}

\subsection*{Problem 74}

Blood phenotype proportions in the U.S.:
\[
P(A)=0.40,\quad P(B)=0.11,\quad P(AB)=0.04,\quad P(O)=0.45.
\]
Assume two randomly selected individuals have independent phenotypes.

\begin{enumerate}
\item[(a)] Probability both phenotypes are $O$:
\[
P(\text{both }O)=P(O)\cdot P(O)=(0.45)^2=0.2025.
\]

\item[(b)] Probability the two phenotypes match:
\[
P(\text{match})=P(A)^2+P(B)^2+P(AB)^2+P(O)^2
=(0.40)^2+(0.11)^2+(0.04)^2+(0.45)^2.
\]
Compute:
\[
0.16+0.0121+0.0016+0.2025=0.3762.
\]
\end{enumerate}

\subsection*{Problem 78}

Let $O_i$ be the event that valve $i$ opens on demand. Given $P(O_i)=0.96$ and the five valves operate
independently.

\begin{enumerate}
\item[(a)] Probability at least one valve opens:
\[
P(\text{at least one opens}) = 1 - P(\text{none open})
= 1 - (0.04)^5
= 1 - 0.0000001024
= 0.9999998976.
\]

\item[(b)] Probability at least one valve fails to open:
Let $F_i$ be the event valve $i$ fails, so $P(F_i)=0.04$.
\[
P(\text{at least one fails}) = 1 - P(\text{none fail})
= 1 - (0.96)^5.
\]
Compute:
\[
0.96^2=0.9216,\quad 0.96^3=0.884736,\quad 0.96^4=0.84934656,\quad 0.96^5=0.8153726976,
\]
so
\[
P(\text{at least one fails}) = 1 - 0.8153726976 = 0.1846273024.
\]
\end{enumerate}

\end{document}
